\documentclass[11pt]{article}
% Shared preamble for the Cupcast papers.
\usepackage[a4paper,margin=2.8cm]{geometry}
\usepackage{amsmath,amssymb,mathtools}
\usepackage{booktabs}
\usepackage{graphicx}
\usepackage{microtype}
% Classic BibTeX via natbib: tectonic runs bibtex natively, while its bundled
% biblatex is version-locked against system biber. Each paper folder symlinks
% references.bib from shared/.
\usepackage[numbers,sort&compress]{natbib}
\usepackage{xcolor}
\usepackage[colorlinks=true,allcolors=blue!45!black]{hyperref}
\graphicspath{{../figures/}}

% Shared notation
\newcommand{\E}{\mathbb{E}}
\newcommand{\Prob}{\mathbb{P}}
\newcommand{\att}{\alpha}
\newcommand{\dfn}{\beta}


\title{Cupcast~II: A Dynamic State-Space Dixon--Coles Model with\\
Squad-Quality Priors and Bounded Context Covariates\\for the 2026 FIFA World Cup}
\author{Carlos Garcia}
\date{June 2026}

\begin{document}
\maketitle

\begin{abstract}
We replace the static Dixon--Coles model of Cupcast~I with a \emph{dynamic
state-space} formulation in which each national team's latent attack and
defence evolve as a Gaussian random walk over quarterly periods, so that form
is a first-class quantity rather than a fixed exponential decay. The initial
strength level of every team is a \emph{learned} function of a lineup-aware
squad-quality composite, built from two seasons of club performance data
normalised by a club-Elo league-strength index and weighted by data-driven
expected national-team minutes. A bounded, audited context layer adjusts
individual match predictions for travel, host crowd, and altitude, each effect
hard-capped so it cannot dominate the statistical core. Inference is by
stochastic variational inference (with No-U-Turn sampling for final fits) in
NumPyro \cite{phan2019numpyro}, and match probabilities feed a 50{,}000-run
Monte Carlo simulation of the full 48-team bracket. Held-out performance against
the actual 2026 group stage, and a comparison against the Cupcast~I baseline,
are reported in Paper~II-V.
\end{abstract}

\section{Introduction}
Cupcast~I \cite{dixoncoles1997} fitted a single static rating per team by
recency-weighted maximum likelihood and shrank it toward an Elo-and-squad prior.
Two weaknesses motivated a rebuild. First, the squad layer was inert: an
attack-only per-90 proxy with thin league coverage left player identity nearly
invisible to the forecast. Second, the model was static and results-driven,
which systematically overweighted recent competitive form (it priced Argentina
as a $16.3\%$ title favourite against a market consensus near $8\%$).

Cupcast~II addresses both directly. Section~\ref{sec:data} describes the data
and the \texttt{clubform} subsystem. Section~\ref{sec:model} specifies the
dynamic state-space goal model and its clubform-informed prior.
Section~\ref{sec:inf} covers inference. Section~\ref{sec:ctx} specifies the
bounded context layer. Section~\ref{sec:sim} describes the tournament
simulation. All stochastic steps are seeded and reproduce exactly.

\section{Data and the Clubform Subsystem}\label{sec:data}
International results, lineups, injuries, and per-fixture player statistics are
pulled from API-Football \cite{apifootball} (all six confederations' World Cup
qualifiers, continental championships, Nations Leagues, and friendlies,
2018--2026; eighteen club leagues over two seasons), cached first and read only
from cache. Youth, B, and non-FIFA selections are filtered out, and matches are
held out at the tournament cutoff.

The \texttt{clubform} subsystem produces, per national team, position-aware
attack/defence/goalkeeping composites. (i)~A \emph{league-strength index}
normalises per-90 statistics across leagues as the median club-Elo
\cite{clubelo} of a league's clubs, mapped to a multiplier centred on one
(Premier League $1.45$ down to weaker leagues near $0.65$), with a curated
fallback for non-European leagues club-Elo does not rate. (ii)~Player quality is
the league-strength-scaled, within-position $z$-score of attacking
(non-penalty goals plus assists, shots, key passes), defensive (tackles,
interceptions, blocks, duels won), and goalkeeping (save rate) output.
(iii)~Expected tournament minutes are estimated from \emph{actual recent
national-team lineups}: a recency-weighted start frequency reconstructs each
nation's core eleven. Per-team composites join player quality to expected
minutes by the global player identifier, with unmatched players imputed at the
position floor; this yields $\approx0.99$ coverage for elite squads.

\section{The Dynamic State-Space Goal Model}\label{sec:model}
For match $m$ at period $p_m$ between home team $h$ and away team $a$, goals are
conditionally Poisson with a Dixon--Coles low-score correction. The log scoring
rates are
\begin{align}
\log\lambda_m &= \mu + \att_h(p_m) - \dfn_a(p_m) + \gamma\,\mathbb{1}[\text{host}_h],\\
\log\nu_m     &= \mu + \att_a(p_m) - \dfn_h(p_m) + \gamma\,\mathbb{1}[\text{host}_a],
\end{align}
with home and away goals $x,y \sim \mathrm{Poisson}(\lambda_m),\mathrm{Poisson}(\nu_m)$
and the joint mass on $\{0,1\}^2$ multiplied by the correction
$\tau(x,y;\lambda,\nu,\rho)$ of \cite{dixoncoles1997}, $\rho\in(-0.3,0.3)$. The
correction is applied as an additive log-likelihood factor.

\paragraph{Dynamics (form).} Each team's attack and defence are time-varying
latent states following a non-centred Gaussian random walk across periods:
\begin{equation}
\att_i(p) = \att_i(0) + \sigma_\att \sum_{t\le p} z^{\att}_{i,t},
\qquad z^{\att}_{i,t}\sim\mathcal N(0,1),
\end{equation}
and analogously for $\dfn_i$ with innovation scale $\sigma_\dfn$. The innovation
scales $\sigma_\att,\sigma_\dfn\sim\mathrm{HalfNormal}$ are the principled
replacement for Cupcast~I's fixed decay half-life; they govern how fast form may
move and, because the filter advances the latent state, in-tournament updating
is automatic. This is the football-specific instance of the dynamic
bivariate-Poisson approach of \cite{koopmanlit2015,baioblangiardo2010}.

\paragraph{Clubform-informed prior $f_\theta$.} The initial level of each team
is a learned linear function of its clubform composite,
\begin{equation}
\att_i(0) \sim \mathcal N\!\big(a_0 + a_1\,\mathrm{clubform}^{\att}_i,\ \tau_{\text{prior}}^2\big),
\qquad a_0,a_1 \sim \mathcal N(0,\cdot),
\end{equation}
with the defensive analogue. Data-rich teams' walks move away from this prior;
data-poor teams remain anchored to it. The learned slope is positive on real
data ($a_1>0$), confirming that club performance predicts international scoring.

\section{Inference}\label{sec:inf}
The model is fitted in NumPyro \cite{phan2019numpyro} on CPU. The primary fitter
is stochastic variational inference with an \texttt{AutoNormal} guide,
initialised at the prior median to avoid the rate overflow a random
initialisation can trigger; No-U-Turn sampling is available for final fits.
Every fit threads an explicit pseudo-random key so results reproduce exactly.

\section{Bounded Context Layer}\label{sec:ctx}
Four of six context families are computed deterministically from the schedule
and a static sixteen-venue table (coordinates, elevation, host nation, and how
late each hosts; Mexican and Canadian venues host only through the Round of~16,
with all later rounds in the United States). Per match and team we compute rest
days, great-circle travel distance from the previous venue, a round-aware host
flag, and venue altitude. Each maps to a hard-capped multiplicative adjustment
of the scoring rate---host $\le e^{0.12}$, travel fatigue $\ge e^{-0.10}$, short
rest $\ge e^{-0.08}$, altitude on the non-adapted side $\ge e^{-0.10}$---with a
per-match provenance record of every raw and applied value. The statistical core
owns the probabilities; the context layer can only nudge. Morale, the one family
without a structured source, is extracted in shadow mode (logged, zero weight)
until an ablation proves its value.

\section{Tournament Simulation}\label{sec:sim}
Match probabilities feed a seeded $50{,}000$-run Monte Carlo simulation of the
official 48-team structure: twelve groups of four under FIFA's group tiebreaker
cascade \cite{fifaregs}, the eight best third-placed teams allocated to the
Round of~32 by FIFA's Annex-C backtracking rule, and knockout rounds through the
final and third-place match. A knockout drawn after ninety minutes proceeds to
extra time, modelled as the same goal model at one-third of the regulation
scoring rates; a tie after extra time is resolved by a penalty shootout held
near even and nudged, within a hard $\pm5$ percentage-point cap, by a
penalty-specific shootout composite built from each squad's goalkeeper
penalty-save and outfield penalty-conversion records (a far more shootout-relevant
signal than general shot-stopping volume, which is confounded by team strength). The simulation runs in about one second
and yields, for every team, the probability of each round, the title, the
runner-up spot, and a podium finish. Proper scoring-rule validation
\cite{gneitingraftery2007} and the Cupcast~I comparison follow in
Paper~II-V.

\bibliographystyle{abbrvnat}
\bibliography{references}

\end{document}
